The goal of this exercise is to write a \texttt{Polynomial} class that represents a polynomial as a set of monomials:

\[
    P(X) = c_0 + c_1 \cdot X + c_2 \cdot X^2 + \dots + c_n \cdot X^n
\]

It uses the \texttt{Monomial} class of Exercise 1, which is provided.

\subsection{Implementation}

Write a \texttt{Polynomial} class with the following requirements:

\begin{itemize}

    \item \textbf{Attributes}:

        \begin{itemize}

            \item \texttt{Monomial[] monomials}: an array where each monomial is stored at the position of its exponent: position 0 holds the monomial with exponent 0, position 1 the one with exponent 1, and so on. If the polynomial has no monomial with a given exponent, that position holds \texttt{null}, which stands for a coefficient of 0. That is why no position may ever hold a monomial whose coefficient is 0.

            \item \texttt{int maxSize}: the length of the array, so the polynomial has room for exponents from 0 to \texttt{maxSize - 1}.

            \item \texttt{int numMonomials}: the number of positions of the array that hold a monomial.

        \end{itemize}

    \item \textbf{Methods}:

        \begin{itemize}

            \item \textbf{Constructors}:

                \begin{itemize}

                    \item A constructor that creates a polynomial with no monomials, given the length of the array.

                    \item A constructor that receives an array of coefficients, where the position of each coefficient is its exponent. A coefficient of 0 gives no monomial.

                    \item A constructor that receives an array of \texttt{Monomial} objects, in any order. The array of the polynomial must have room for the highest exponent among them.

                \end{itemize}

            \item \texttt{getMonomial()}: returns the monomial with a given exponent, or \texttt{null} if there is none or the exponent is outside the array.

            \item \texttt{addMonomial()}: adds a monomial to the polynomial, at the position of its exponent. If there is already a monomial there, it is replaced by the sum of both. If the coefficient of the result is 0, or within the tolerance of \texttt{isEqualTo()} of it, such as when \(3X^2\) and \(-3X^2\) are added, the position is left empty, with \texttt{null}. If the array has no room for that exponent, the polynomial does not change.

            \item \texttt{replaceMonomial()}: replaces the monomial that has the same exponent as a given one, or removes it if the coefficient of the given one is 0. If the polynomial has no monomial with that exponent, it does not change.

            \item \texttt{add()}: adds another polynomial to this one and returns a \textbf{new} \texttt{Polynomial}.

            \item \texttt{scalarProduct()}: multiplies the polynomial by a number and returns a \textbf{new} \texttt{Polynomial}.

            \item \texttt{product()}: multiplies this polynomial by another one and returns a \textbf{new} \texttt{Polynomial}: every monomial of the first is multiplied by every monomial of the second, and the products are added.

            \item \texttt{solve()}: given a value for \(X\), returns the value of the polynomial.

            \item \texttt{describe()}: returns a \texttt{String} with the polynomial, such as \texttt{"1.0X\textasciicircum{}0 + 3.0X\textasciicircum{}2"}, or \texttt{"0"} if it has no monomials. \texttt{numMonomials} tells you whether that is the case. The \texttt{describe()} method of \texttt{Monomial} gives the text of each monomial.

            \item \texttt{isEqualTo()}: given another polynomial, checks whether both are equal, that is, whether they have the same coefficient for every exponent. A missing monomial counts as a coefficient of 0, so two equal polynomials may have arrays of different lengths. As in Exercise 4, compare the coefficients with a small tolerance.

        \end{itemize}

\end{itemize}

None of the operations changes the polynomials it works with.

\subsection{Recommendations}

\begin{itemize}

    \item \textbf{Incremental implementation}: write and test each method on its own before moving on to the next one, making sure each part works before building the rest of the class.

    \item \textbf{Reuse}: \texttt{addMonomial()} already adds a monomial to the one in its position and leaves out the ones that cancel out, so the constructors, \texttt{add()}, \texttt{scalarProduct()} and \texttt{product()} become simpler if they use it. Call it on the polynomial being built, as \texttt{result.addMonomial(...)}, or on the object itself, as \texttt{this.addMonomial(...)}.

    \item \textbf{Building the text}: \texttt{describe()} joins one piece of text per monomial, so a \texttt{StringBuilder} is a better fit than \texttt{+}.

    \item \textbf{Keeping} \texttt{numMonomials} \textbf{right}: every time a position goes from \texttt{null} to a monomial, or back, update it.

\end{itemize}

\subsection{Evaluation and testing}

To complete the exercise, write a \texttt{Main} class with a \texttt{main} method that performs the following tests:

\begin{enumerate}

    \item \textbf{Object creation}:

        \begin{itemize}

            \item Create polynomials with each of the three constructors.

        \end{itemize}

    \item \textbf{Operations and verification}:

        \begin{itemize}

            \item Operate with polynomials and check that the results are the expected ones, including a sum whose coefficients cancel out.

            \item Evaluate polynomials for some values of \(X\) to check the results.

            \item Check whether different polynomials are equal, including the same polynomial built in two different ways.

        \end{itemize}

    \item \textbf{Clean code and documentation}:

        \begin{itemize}

            \item Make sure the code is clean and legible, with explanatory comments where necessary.

        \end{itemize}

\end{enumerate}
